\section{Proof of an Implication}

\begin{frame}{Proof of an Implication}
\begin{center}
\vspace{1.5cm}
{\Huge\textbf{Proof of an Implication}}
\end{center}
\end{frame}

% --------------------------------------------------
% Question 1
% --------------------------------------------------

\begin{frame}{Question 1}
A software system claims that if a user's access score is divisible by $4$,
then the score is even.

Prove the implication directly for every integer access score.
\end{frame}

\begin{frame}{Question 1 -- Answer}
Let $n$ be an integer such that:

$$
4\mid n
$$

Then, by the definition of divisibility, there exists an integer $k$ such
that:

$$
n=4k
$$

\end{frame}

\begin{frame}{Question 1 -- Answer}
Rewrite $n$ as:

$$
n=2(2k)
$$

Since $2k$ is an integer, $n$ has the form $2m$ for an integer $m$.

Therefore, $n$ is even.

Hence:

$$
\boxed{4\mid n\Rightarrow2\mid n}
$$

\end{frame}

% --------------------------------------------------
% Question 2
% --------------------------------------------------

\begin{frame}{Question 2}
A cloud security system follows the rule:

\begin{center}
\textit{If a server passes authentication, then it is allowed to access the network.}
\end{center}

A particular server has passed authentication.

\textbf{Question:} Using Modus Ponens, determine whether the server is
allowed to access the network.
\end{frame}

\begin{frame}{Question 2 -- Answer}
Let:

$$
P=\text{The server passes authentication}
$$

and

$$
Q=\text{The server is allowed to access the network}
$$

The security rule is:

$$
P\rightarrow Q
$$

The given information states:

$$
P
$$

\end{frame}

\begin{frame}{Question 2 -- Answer}
We therefore have:

$$
P\rightarrow Q
$$

and

$$
P
$$

Modus Ponens states:

$$
P\rightarrow Q,\quad P
\quad\therefore\quad Q
$$

Therefore:

$$
\boxed{Q}
$$

The server is allowed to access the network.
\end{frame}

% --------------------------------------------------
% Question 3
% --------------------------------------------------

\begin{frame}{Question 3}
A machine-learning deployment system follows the rule:

\begin{center}
\textit{If a model is deployed, then it has passed all required tests.}
\end{center}

An administrator verifies that the model has \textbf{not} passed all the
required tests.

\textbf{Question:} Using Modus Tollens, determine whether the model can be
considered deployed under this rule.
\end{frame}

\begin{frame}{Question 3 -- Answer}
Let:

$$
P=\text{The model is deployed}
$$

and

$$
Q=\text{The model has passed all required tests}
$$

The rule is:

$$
P\rightarrow Q
$$

The given information is:

$$
\neg Q
$$

\end{frame}

\begin{frame}{Question 3 -- Answer}
Modus Tollens states:

$$
P\rightarrow Q,\quad\neg Q
\quad\therefore\quad\neg P
$$

Since the model has not passed all required tests,

$$
\neg Q
$$

Therefore:

$$
\boxed{\neg P}
$$

The model cannot be considered deployed under the given rule.
\end{frame}

% --------------------------------------------------
% Question 4
% --------------------------------------------------

\begin{frame}{Question 4}
A program produces the value:

$$
f(n)=n^2+n
$$

for an integer input $n$.

Prove that $f(n)$ is even for every integer $n$.
\end{frame}

\begin{frame}{Question 4 -- Answer}
We need to prove:

$$
n\in\mathbb Z\Rightarrow n^2+n\text{ is even}
$$

Factor the expression:

$$
n^2+n=n(n+1)
$$

\end{frame}

\begin{frame}{Question 4 -- Answer}
Among any two consecutive integers, one must be even.

Therefore, either:

$$
2\mid n
$$

or:

$$
2\mid(n+1)
$$

Thus, the product

$$
n(n+1)
$$

must be even.
\end{frame}

\begin{frame}{Question 4 -- Answer}
Hence:

$$
\boxed{n^2+n\text{ is even}}
$$

for every integer $n$.
\end{frame}

% --------------------------------------------------
% Question 5
% --------------------------------------------------

\begin{frame}{Question 5}
A computing system represents a positive integer $n$ using its binary
representation.

Prove that if $n$ is odd, then its binary representation must end in $1$.
\end{frame}

\begin{frame}{Question 5 -- Answer}
Every positive integer can be written as:

$$
n=2q+r
$$

where:

$$
r\in\{0,1\}
$$

The remainder $r$ determines the last binary digit.
\end{frame}

\begin{frame}{Question 5 -- Answer}
If $n$ is odd, then division by $2$ gives remainder $1$.

Therefore:

$$
n=2q+1
$$

Hence:

$$
r=1
$$

\end{frame}

\begin{frame}{Question 5 -- Answer}
In binary representation, the remainder obtained after division by $2$
corresponds to the least significant bit.

Therefore:

$$
\boxed{\text{An odd integer has a binary representation ending in }1}
$$

\end{frame}
